\begin{abstract}
Evaluating large language model (LLM) based chat assistants is challenging due to their broad capabilities and the inadequacy of existing benchmarks in measuring human preferences.
To address this, we explore using strong LLMs as judges to evaluate these models on more open-ended questions.
We examine the usage and limitations of LLM-as-a-judge, including position, verbosity, and self-enhancement biases, as well as limited reasoning ability, and propose solutions to mitigate some of them.
We then verify the agreement between LLM judges and human preferences by introducing two benchmarks: MT-bench, a multi-turn question set; and Chatbot Arena, a crowdsourced battle platform.
Our results reveal that strong LLM judges like GPT-4 can match both controlled and crowdsourced human preferences well, achieving over 80\% agreement, the same level of agreement between humans.
Hence, LLM-as-a-judge is a scalable and explainable way to approximate human preferences, which are otherwise very expensive to obtain.
Additionally, we show our benchmark and traditional benchmarks complement each other by evaluating several variants of LLaMA and Vicuna.
The MT-bench questions, 3K expert votes, and 30K conversations with human preferences are publicly available at \url{https://github.com/lm-sys/FastChat/tree/main/fastchat/llm\_judge}.
\end{abstract}

\section{Introduction}
\label{sec:introduction}
There has been a proliferation of LLM-based chat assistants (chatbots) that leverage supervised instruction fine-tuning and reinforcement learning with human feedback (RLHF) to unlock new instruction following and conversational abilities~\cite{ouyang2022training,bai2022training,openai2023gpt4,vicuna2023,zhou2023lima,xu2023wizardlm,dubois2023alpacafarm}.
Once aligned with humans, these chat models are strongly preferred by human users over the original, unaligned models on which they are built.
However, the heightened user preference does not always correspond to improved scores on traditional LLM benchmarks -- benchmarks like MMLU~\cite{hendrycks2020measuring} and HELM~\cite{liang2022holistic} cannot effectively tell the difference between these aligned models and the base models.
This phenomenon suggests that there is a fundamental discrepancy between user perceptions of the usefulness of chatbots and the criteria adopted by conventional benchmarks.

We argue that this discrepancy primarily arises due to  
existing evaluation that only measures LLMs' core capability on a confined set of tasks (e.g., multi-choice knowledge or retrieval questions), without adequately assessing its alignment with human preference in open-ended tasks, such as the ability to accurately adhere to instructions in multi-turn dialogues.
As a demonstration, we show conversation histories with two models on an MMLU question in \Cref{fig:intro-figure}.
The two models are LLaMA-13B~\cite{touvron2023llama}, a pre-trained base model without fine-tuning, and Vicuna-13B, our fine-tuned model from LLaMA-13B on high-quality conversations (the training details are in \Cref{sec:training-details}).
Despite the base LLaMA models showing competitive performance on conventional benchmarks (\Cref{table:dataset_composition}), its answers to open-ended questions are often not preferred by humans.
This misalignment of conventional benchmarks underscores the core problem driving this paper:
\emph{the need for a robust and scalable automated method to evaluate LLM alignment with human preferences.}

\begin{figure}[t]
    \centering
    \vspace{-1em}
    \includegraphics[width=0.95\linewidth]{figures/intro_fed.pdf}
    \vspace{-0.5em}
    \caption{Multi-turn dialogues between a user and two AI assistants—LLaMA-13B (Assistant A) and Vicuna-13B (Assistant B)—initiated by a question from the MMLU benchmark and a follow-up instruction. GPT-4 is then presented with the context to determine which assistant answers better.}
    \label{fig:intro-figure}
    \vspace{-1.5em}
\end{figure}

To study this, we introduce two benchmarks with human ratings as the primary evaluation metric: MT-bench and Chatbot Arena.
MT-bench is a series of open-ended questions that evaluate a chatbot's multi-turn conversational and instruction-following ability -- two critical elements for human preference. MT-bench is also carefully constructed to differentiate chatbots based on their core capabilities, such as reasoning and math.
In addition, we develop Chatbot Arena, a crowdsourced platform featuring anonymous battles between chatbots in real-world scenarios -- Users engage in conversations with two chatbots at the same time and rate their responses based on personal preferences. 



While human evaluation is the gold standard for assessing human preferences, it is exceptionally slow and costly.
To automate the evaluation, we explore the use of state-of-the-art LLMs, such as GPT-4, as a surrogate for humans. 
Because these models are often trained with RLHF, they already exhibit strong human alignment.
We call this approach \emph{``LLM-as-a-judge''}. 
This approach has been tried in our earlier blog post~\cite{vicuna2023} and other concurrent or follow-up work~\cite{bubeck2023sparks, openaievals, dubois2023alpacafarm, dettmers2023qlora, zhou2023lima,gudibande2023false,peng2023instruction,wang2023large,chiang2023can,wang2023far}.
However, there has not been a systematic study of this approach. %


In this paper, we study the LLM-as-a-judge approach by comparing it to the gold standard of human evaluation. 
We examine several potential limitations of the LLM-as-a-judge approach including position bias, verbosity bias, self-enhancement bias, and limited reasoning ability.
We show that some of the biases are minor or can be mitigated. 
Once addressed, our results from 3K controlled expert votes and 3K crowdsourced human votes in the wild verify that GPT-4 judge match human evaluations at an agreement rate exceeding 80\%, achieving the same level of human-human agreement (\S\ref{sec:agreement}, Table~\ref{tab:failure_rate_math}).
Consequently, this suggests LLM-as-a-judge is a scalable method to swiftly evaluate human preference, serving as a promising alternative to traditional human evaluations.  \\

This paper makes two contributions: (1) a systematic study of LLM-as-a-judge; and (2) human preference datasets with high-quality questions and diverse user interactions from MT-bench and Chatbot Arena.
In addition, we argue for the adoption of a hybrid evaluation framework for future LLM benchmarks: by combining the existing capability-based benchmarks and the new preference-based benchmarks with LLM-as-a-judge, one can swiftly and automatically evaluate both the core capabilities and human alignment of models.
We publicly release 80 MT-bench questions, 3K expert votes, and 30K conversations with human preferences for future study.



























